\subsection*{Related Work}
\paragraph{Adaptive Updates and Convergence.}
Adaptive coordinate scaling predates Adam through AdaGrad and ADADELTA
\cite{duchi2011adagrad,zeiler2012adadelta}. Adam uses gradient-moment
averages \cite{kingma2015adam}; AdamW separates weight decay from that
scaling \cite{loshchilov2019adamw}. Sign and variance interpretations
\cite{balles2018dissecting,bernstein2018signsgd} and alternative
denominators such as AdaBelief \cite{zhuang2020adabelief} are established.
Convergence analyses cover nonconvex settings
\cite{zaheer2018adaptive,chen2019adamtype}, averaging-parameter conditions
\cite{zou2019sufficient}, and unchanged Adam updates
\cite{defossez2022simple,zhang2022adam}. Their assumptions do not imply
minimizer consistency for every fixed-memory, constant-step problem.
Our result concerns the specified AdamW denominator, including its shared
history with the gradient.

\paragraph{The Closest Comparisons.}
The averaged Adam field already depends on the whole sampling law
\cite{dj2024}. Its derivative and corrected noise covariance enter a
decreasing-step, averaged-iterate central limit theorem conditional on
convergence to a stable zero \cite{dj2026}. Theorem~4.12 of
\cite{symmetry} concerns fixed isotropic curvature; here sign-symmetric
noise and random curvature preserve the same zero while changing stability.
Stochastic convex failure \cite{reddi}, deterministic AdamW instability
and exact discrete-step conditions
\cite[Proposition~4.3 and Appendix~A.6]{su}, and large-step stability
phenomena \cite{cohen_edge,adaptive_edge,eos} are prior results.
Normalization-dependent floors versus SGD under multiplicative noise also
precede us \cite{singh}; that comparison changes the normalization rule.
Our distinction is the same AdamW update, with fixed positive offset and
decay, having vanishing versus persistent stationary error when only the
pairing changes. The average loss, full gradient covariance kernel and both
complete Hessian and noise marginals remain fixed. The moment averages
continue fluctuating at the minimum even as the step vanishes.

\paragraph{Classical Methods and Stationary Error.}
Stochastic approximation \cite{robbinsmonro1951}, its differential-equation
and stability methods \cite{kushneryin2003,benaim1999,borkarmeyn2000},
unstable-point avoidance \cite{pemantle1990}, and constant-step stationary
limits \cite{fortpages,ddb,chenmoumaguluri2022} supply established tools.
We prove their required gradient-history estimates rather than claiming new
Gaussian-limit or covariance methods. The additional work controls every
stationary distribution satisfying the algorithm's moment constraints,
including a case where
the optimum repels the averaged dynamics. Exact SGD balance identities
\cite{yaida2019} presuppose stationarity; martingale corrections handle
dependent noise \cite{fort2015}. Iterate averaging changes fluctuations
\cite{polyakjuditsky1992,mouetal2020}; our error concerns an unaveraged
stationary iterate. Classical rearrangement inequalities
\cite{cambanissimonsstout1976} give our scalar extrema after deriving the
optimizer-specific response weight, not a new general coupling inequality.

\paragraph{Diffusion Models and Memory.}
Finite-time stochastic differential equation approximations
\cite{litaie2017} extend to adaptive methods and scaling rules
\cite{malladi,compagnoni}. Their Adam limits let memory parameters approach
one as the step decreases; our main theorem fixes memory and the sampling
law. These approximations already describe interactions between adaptive
normalization and noise; different memory scalings can retain different
parts of that interaction. Finite-time approximation alone does not justify
exchanging long-time and small-step limits. Our separate long-memory result takes the small-step
limit first, without covering every simultaneous scaling. The example
challenges descriptions retaining only the matched summaries, not diffusion
models that retain additional dependence information.

\paragraph{Noise, Curvature and Adaptive Bias.}
SGD diffusion models connect fluctuations to approximate Bayesian inference
and noise scale \cite{mandt2017,smithle2018}, and anisotropic noise to
probability currents \cite{chaudharisoatto2018}. Curvature--noise tradeoffs
\cite{thomas}, Adam covariance with a population replacement for the
random denominator \cite{liu2021}, bias near minimizer manifolds under
another memory scaling \cite{liwenlyu}, preconditioning limits \cite{das},
and discrete-update regularization \cite{cattaneo} are established.
Keeping the realized denominator random distinguishes our exact response
formula; the matched-law stationary error orders distinguish the
construction.

\paragraph{Connections to Transformers and Limits.}
Heavy-tailed gradient noise \cite{simsekli,zhang_attention} and Hessian
block differences \cite{zhang_hessian} offer other explanations for
adaptive optimization. SGD--Adam gaps can persist without sampling noise
\cite{kunstner2023noise}, and optimization gains need not improve
generalization \cite{wilson2017marginal}. Our bounded-noise examples fix
the entire Hessian law and isolate an additional dependence. They establish
neither a full explanation of transformer optimization nor a language-model
or generalization benefit. Bounded noise makes heavy tails unnecessary for
this particular separation, without disputing their empirical role in
neural-network training.

